\documentclass[11pt,a4paper]{article}
\usepackage[margin=26mm]{geometry}
\usepackage[T1]{fontenc}
\usepackage{lmodern}
\usepackage{amsmath,amssymb,amsthm}
\usepackage{microtype}
\usepackage[hidelinks]{hyperref}
\newtheorem{theorem}{Theorem}
\newtheorem{lemma}[theorem]{Lemma}
\newtheorem{corollary}[theorem]{Corollary}
\theoremstyle{definition}
\newtheorem{definition}[theorem]{Definition}
\newcommand{\Ad}{\operatorname{Ad}}
\newcommand{\U}{\mathcal U}
\newcommand{\T}{\mathbb T}
\newcommand{\CK}{C^*}
\title{Strong approximate innerness of finite abelian\protect\\
diagonal actions on Cuntz--Krieger algebras}
\author{Alec Kriebel\\\small Independent researcher\\
\small\href{https://orcid.org/0009-0001-9320-500X}{ORCID: 0009-0001-9320-500X}}
\date{17 October 2026}
\begin{document}
\maketitle
\begin{abstract}
We give a complete affirmative answer to the historical Question~1 in
Barlak's Oberwolfach report on diagonal quasi-free automorphisms: every
finite abelian diagonal quasi-free action on a simple purely infinite
Cuntz--Krieger algebra associated with a finite essential $0$--$1$ matrix
is strongly approximately inner. The proof is an explicit deduction from
Gabe--Szab\'o's equivariant uniqueness theorem, together with the established
diagonal-coboundary mechanism and elementary finite-group averaging.
It covers periodic matrices, nonouter actions and nonfaithful actions.
This note is an exposition and verification of consequences of prior
theorems; it makes no claim of a first solution or a new mechanism.
\end{abstract}

\noindent\textbf{AI-use and verification disclosure.}
Extensive AI assistance was used in mathematical reasoning, proof drafting,
source auditing, diagnostic computations and manuscript review. This
manuscript is unrefereed; no external human peer review or formal
proof-assistant certification is claimed. Finite computations in the
accompanying verification material are diagnostics, not proofs of the
operator-algebra assertions.

\section{Statement and prior results}

An action $\alpha:G\curvearrowright A$ on a unital $C^*$-algebra is
\emph{strongly approximately inner} if, for every $g\in G$, there are
$u_{g,n}\in\U(A^\alpha)$ such that
\[
 \lim_n\|u_{g,n}a u_{g,n}^*-\alpha_g(a)\|=0\qquad(a\in A).
\]
Here $A^\alpha$ is fixed by the \emph{entire} action. The sequences for
different $g$ need not satisfy any approximate representation relation.
This is the property in Izumi's definition and in
\cite[Definition~2.5]{BS}; it is weaker than approximate representability.

Let $M$ be a finite $0$--$1$ matrix with no zero rows or columns. Write
$\mathcal O_M$ for its Cuntz--Krieger algebra, with canonical generators
$s_i$ and mutually orthogonal range projections $p_i=s_i s_i^*$ summing
to $1$, where $s_i^*s_i=\sum_j M(i,j)p_j$. We use the original workshop
definition of diagonal quasi-free: the action preserves
$\operatorname{span}\{s_i\}$ and fixes every $p_i$.

\begin{theorem}\label{thm:ck}
If $\mathcal O_M$ is simple and purely infinite, every action of a finite
abelian group on $\mathcal O_M$ by these diagonal quasi-free automorphisms
is strongly approximately inner.
\end{theorem}

This answers the full question in \cite[p.~2621, Question~1]{OWR}, without
the outerness and aperiodicity assumptions of the adjacent positive
result. The report has bibliographic year 2017 and was published online
on 5 July 2018. Izumi had already established the standard finite Cuntz
algebra case without outerness; the precise inspected source is his
2002 preprint, Proposition~5.6(2) \cite{Izumi}. Barlak--Szab\'o proved the
outer aperiodic Cuntz--Krieger and graph cases
\cite[Theorems~1--2]{BS}; these are Corollaries~2.7--2.8 in their arXiv v3.

Our shorter proof uses the general uniqueness theorem of Gabe--Szab\'o,
first available in 2022 \cite{GS}. It does not assert that they printed
Theorem~\ref{thm:ck}. The diagonal-inner argument below is the finite-graph
form of Matsumoto's earlier coboundary and periodic-point mechanism
\cite[Lemma~5.14 and Proposition~6.4]{Matsumoto}. An alternative explicit
return-corner proof from the Barlak--Szab\'o theorem is retained in the
verification material. Its graph corner is a special case of
Crisp--Gow \cite[Theorem~3.1 in v1]{CG}. These are prior mechanisms.
The present source audit does not establish priority for an explicit
older statement of the combined full theorem. In particular, the full
Zacharias chapter and Izumi's final 2004 typeset body were not available
in that audit; negative search results give no priority certificate.

\section{Two finite-group facts}

We first record the exact fixed-algebra correction that will be used twice.
If $x_n$ are unitaries and
$\max_{h\in G}\|\gamma_h(x_n)-x_n\|\to0$ for a finite-group action
$\gamma$, then
\[
 b_n=|G|^{-1}\sum_{h\in G}\gamma_h(x_n)\in A^\gamma,
 \qquad \|b_n-x_n\|\longrightarrow0.
\]
Thus $b_n$ is eventually invertible and its polar unitary
$b_n(b_n^*b_n)^{-1/2}$ belongs to $A^\gamma$ and is norm-close to $x_n$.
Replacing $x_n$ by this unitary preserves any point-norm implementing limit.

\begin{lemma}\label{lem:cocycle}
Let $G$ be finite abelian and $\beta:G\curvearrowright A$ strongly
approximately inner. If $d:G\to\U(A)$ is an honest $\beta$-cocycle,
so $d_{gh}=d_g\beta_g(d_h)$ and $d_1=1$, then
$\alpha_g=\Ad(d_g)\beta_g$ is strongly approximately inner.
\end{lemma}
\begin{proof}
Fix $g$ and choose $v_n\in\U(A^\beta)$ implementing $\beta_g$ in the
point-norm limit. Put $t_n=d_gv_n$. These unitaries implement $\alpha_g$
asymptotically. For $h\in G$, abelianness and the cocycle law give
\[
 \alpha_h(t_n)=d_h\beta_h(d_g)v_nd_h^*
              =d_g\beta_g(d_h)v_nd_h^*.
\]
Consequently
\[
 \|\alpha_h(t_n)-t_n\|
   =\|\beta_g(d_h)-v_nd_hv_n^*\|\longrightarrow0.
\]
There are only finitely many $h$. Averaging over $\alpha$ and taking the
polar unitary as above produces unitaries in $A^\alpha$ with the same
implementing limit.
\end{proof}

The hypothesis is an actual cocycle law, not a scalar-projective law for
arbitrary pointwise inner implementers. It imposes no group law on the
final approximating sequences. The converse follows by applying the lemma
to the honest $\alpha$-cocycle $g\mapsto d_g^*$.
The general framework of \cite[Remark~3.2 and Theorem~3.7]{CNS} supplies
prior duality results for precisely this weaker strong approximate
innerness property; the displayed elementary proof makes the transfer
needed here explicit.

\begin{lemma}\label{lem:outer}
Let $G$ be finite abelian and let $\gamma:G\curvearrowright A$ be a
pointwise outer action on a unital Kirchberg algebra. Suppose every
$\gamma_g$ is homotopic to $\mathrm{id}_A$ through unital
$G$-equivariant homomorphisms. Then $\gamma$ is strongly approximately inner.
\end{lemma}
\begin{proof}
Apply Gabe--Szab\'o's unital equivariant uniqueness theorem
\cite[Theorem~5.8]{GS} to the cocycle embeddings
$(\mathrm{id}_A,1)$ and $(\gamma_g,1)$, separately for each $g$.
Here $G$ is exact; $A$ is separable, unital and exact; the finite-group
action is amenable; and pointwise outerness implies isometric shift
absorption by \cite[Proposition~3.14 in v1]{GS}.
The proposed homotopy is equivariant, so the two embeddings have equal
$KK^G$ classes. Theorem~5.8 and Definition~1.15 of that source give
unitaries $u_t\in\U(A)$ with
\[
 \Ad(u_t)(a)\longrightarrow\gamma_g(a),\qquad
 \max_{h\in G}\|\gamma_h(u_t)-u_t\|\longrightarrow0.
\]
The second conclusion is the cocycle-defect limit for the two trivial
cocycles. The target is unital, so these are unitaries in $A$ itself.
Average and take polar unitaries, then sample times tending to infinity.
This gives the required sequence in $A^\gamma$.
\end{proof}

For clarity, the cited v1 locators are Definition~1.15, p.~13;
Proposition~3.14, p.~29; and Theorem~5.8, pp.~45--46. In v3 the proposition
is numbered 3.15 and Theorem~5.8 is on p.~46. We use these inspected
preprint statements, with their version-specific numbers, rather than
guessing final article numbers. The unital theorem has no UCT hypothesis,
additional stable-target hypothesis, or requirement that the implementing
path start at $1$.

\section{Edge phases and the inner kernel}

Let $E$ be a finite strongly connected directed graph with at least one
edge which is not a directed cycle. Use graph generators $p_v,S_e$ with
\[
 S_e^*S_e=p_{r(e)},\qquad
 p_v=\sum_{s(e)=v}S_eS_e^*.
\]
Distinct edges have orthogonal range projections. For a path $\mu$, use
$S_\mu$ for its product, and include vertices as paths of length zero.
An edge-diagonal action has the form
\[
 \alpha_g(p_v)=p_v,\qquad \alpha_g(S_e)=\chi_e(g)S_e,
\]
where $\chi_e:G\to\T$ is a character. Set
$\chi_\mu=\prod_{e\text{ in }\mu}\chi_e$, with vertex character $1$.

\begin{lemma}\label{lem:inner}
If $\sigma(S_e)=z_eS_e$ and $\sigma(p_v)=p_v$, then $\sigma$ is inner
if and only if there are $t_v\in\T$ with
$z_e=t_{s(e)}\overline{t_{r(e)}}$. In that case
$\sigma=\Ad(\sum_vt_vp_v)$. The $t_v$ are unique after normalizing one of
them to $1$.
\end{lemma}
\begin{proof}
Every cycle has an exit: otherwise strong connectivity would make that
exitless cycle the whole graph. The path diagonal
$D_E=\overline{\operatorname{span}}\{S_\mu S_\mu^*\}$ is therefore a
MASA \cite[Theorem~5.2]{HPP}; see also the convention-compatible statement
in \cite[\S2.1]{CHS}. Since $\sigma$ fixes $D_E$ pointwise, an inner
implementer $h$ must lie in $\U(D_E)$.

Identify $D_E$ with $C(E^\infty)$, with path-range projections the cylinder
functions. The relation $hS_\mu=z_\mu S_\mu h$ gives
$h(\mu x)=z_\mu h(x)$ for every infinite tail $x$ starting at $r(\mu)$.
For a closed path $\mu$ choose $x=\mu\mu\cdots$. Then $\mu x=x$ and
$|h(x)|=1$, so $z_\mu=1$.

Fix a base vertex $b$ and directed paths $\theta_v:b\to v$, with
$\theta_b$ empty. Put $t_v=z_{\theta_v}^{-1}$. If $e:v\to w$, append
any path $\eta:w\to b$. Both $\theta_v e\eta$ and $\theta_w\eta$ are
closed, giving $z_{\theta_v}z_e=z_{\theta_w}$ and hence the claimed
formula. Conversely that formula makes conjugation by $\sum_vt_vp_v$
agree with $\sigma$ on generators. Two such potentials have ratio
constant along edges, hence constant by connectivity; normalization gives
uniqueness.
\end{proof}

\begin{theorem}\label{thm:graph}
Every edge-diagonal action of a finite abelian group on $C^*(E)$, for $E$
as above, is strongly approximately inner.
\end{theorem}
\begin{proof}
Put $A=C^*(E)$, a unital Kirchberg algebra by the standard graph
simplicity and pure-infiniteness criteria. First suppose the action is
pointwise outer. For fixed $g$, choose real $a_e$ with
$e^{ia_e}=\chi_e(g)$. The assignments
\[
 \delta_t(p_v)=p_v,\qquad
 \delta_t(S_e)=e^{ita_e}S_e\quad(0\leq t\leq1)
\]
preserve the graph relations and have inverse phases, so give
automorphisms. They are point-norm continuous, first on generator
polynomials and then by density. All commute with every $\alpha_h$.
They form an equivariant homotopy from $\mathrm{id}$ to $\alpha_g$.
Lemma~\ref{lem:outer} applies. No graph-period condition occurs.

For the general action let
$K=\{k\in G:\alpha_k\text{ is inner}\}$, the kernel of the homomorphism
$G\to\operatorname{Out}(A)$. Fix the paths $\theta_v$ in
Lemma~\ref{lem:inner}, once for all group elements, and define on all of $G$
\[
 D_g=\sum_v\chi_{\theta_v}(g)^{-1}p_v.
\]
These unitaries form a genuine representation: $D_gD_h=D_{gh}$.
They are fixed by every $\alpha_h$. Lemma~\ref{lem:inner} gives
$\alpha_k=\Ad(D_k)$ for $k\in K$. Thus
\[
 \beta_g=\Ad(D_g^*)\alpha_g
\]
is an edge-diagonal action, trivial on $K$, with phases
$\chi_{\theta_{s(e)}}(g)\chi_e(g)\chi_{\theta_{r(e)}}(g)^{-1}$.
It descends to $H=G/K$. This quotient action is pointwise outer:
if $\beta_g$ is inner, then $\alpha_g=\Ad(D_g)\beta_g$ is inner, hence
$g\in K$. The outer case proves strong approximate innerness of the
$H$-action and therefore of its inflation $\beta$, with the same fixed
algebra. If $H$ is trivial, $\beta$ is trivial and the unit implements
every $\beta_g$ exactly.

Finally $\beta_g(D_h)=D_h$ and $D_gD_h=D_{gh}$, so $D$ is an honest
$\beta$-cocycle and $\alpha_g=\Ad(D_g)\beta_g$.
Lemma~\ref{lem:cocycle} proves the assertion with unitaries in the
original $A^\alpha$. This also covers $K=G$, where the $D_g$ already
implement $\alpha_g$ exactly in $A^\alpha$, and covers nonfaithful actions.
\end{proof}

\section{The original matrix presentation}

\begin{proof}[Proof of Theorem~\ref{thm:ck}]
Write $\alpha_g(s_i)=\sum_j c_{ji}(g)s_j$. For $j\ne i$, multiply by
$p_j$: the fixed range condition gives
$p_j\alpha_g(s_i)=0$, hence $c_{ji}(g)s_j=0$. Thus
$\alpha_g(s_i)=\chi_i(g)s_i$, where $\chi_i$ is a character by the action
law. No simultaneous diagonalization is needed for this original
definition. The broader domain-projection convention in \cite{BS} can
instead be reduced to eigenvectors using its Proposition~1.8(2).

Let $E$ have vertices $1,\ldots,n$ and an edge $i\to j$ whenever
$M(i,j)=1$. In $\mathcal O_M$ put $T_{ij}=s_ip_j$ for these edges.
Then
\[
 T_{ij}^*T_{ij}=p_j,\qquad
 \sum_{j:M(i,j)=1}T_{ij}T_{ij}^*=p_i.
\]
These give a graph Cuntz--Krieger family. Conversely the outgoing sums
$s_i=\sum_{j:M(i,j)=1}T_{ij}$ in $C^*(E)$ satisfy the matrix relations;
cross terms vanish since the edge initial projections $p_j$ are
orthogonal. The two formulas are inverse on generators. Thus
$\mathcal O_M\cong C^*(E)$ equivariantly, with edge phases $\chi_i$.

Essentiality gives at least one edge and both no sinks and no sources. In this finite setting,
simplicity forces strong connectivity. Indeed, if the condensation graph
has more than one component, choose a terminal component $T$ and a
different source component $S$. Since the original graph has no sources,
$S$ contains a cycle. The hereditary set $T$ is nonempty and its saturated
closure cannot absorb that cycle: at each stage, every still excluded
cycle vertex has its next cycle edge outside the set. Its saturated
closure is therefore proper. The associated nonzero proper graph ideal
contradicts simplicity. This uses the standard saturated-hereditary ideal
construction for graph algebras \cite{KPR}.

A directed cycle has algebra $M_m(C(\T))$, so is excluded by the
hypotheses. Theorem~\ref{thm:graph} applies and proves the full assertion.
\end{proof}

\noindent\textbf{Scope.}
All sequences obtained above lie exactly in the fixed algebra of the
whole original action. The proof does not replace this condition by
ordinary approximate innerness or by fixedness under a single group
element. It uses neither matrix primitivity nor action faithfulness.
The verification package preserves the longer return-corner deduction
and the versioned source checks separately from this note; neither
AI review nor those diagnostics certify novelty or replace the cited
mathematical inputs.

\begin{thebibliography}{10}\small
\bibitem{OWR} S.~Barlak, Diagonal quasi-free automorphisms of
Cuntz--Krieger algebras, in \emph{Oberwolfach Reports} \textbf{14}
(2017), no.~3, report 42, p.~2621.
\href{https://doi.org/10.4171/OWR/2017/42}{doi:10.4171/OWR/2017/42}.

\bibitem{BS} S.~Barlak and G.~Szab\'o, On diagonal quasi-free
automorphisms of simple Cuntz--Krieger algebras,
\emph{Math. Scand.} \textbf{125} (2019), 210--226.
\href{https://doi.org/10.7146/math.scand.a-114823}{doi:10.7146/math.scand.a-114823}.
Versioned source: \href{https://arxiv.org/abs/1708.08258v3}{arXiv:1708.08258v3}.

\bibitem{CHS} R.~Conti, J.~H.~Hong and W.~Szyma\'nski, Endomorphisms
of graph algebras, \emph{J. Funct. Anal.} \textbf{263} (2012), 2529--2554.
\href{https://doi.org/10.1016/j.jfa.2012.08.024}{doi:10.1016/j.jfa.2012.08.024}.
Versioned source: \href{https://arxiv.org/abs/1101.4210v3}{arXiv:1101.4210v3}, \S2.1, p.~5.

\bibitem{CNS} J.~Christensen, R.~Neagu and G.~Szab\'o, On a Rokhlin
property for abelian group actions on $C^*$-algebras,
\emph{J. Noncommut. Geom.}, published online 10 December 2025.
\href{https://doi.org/10.4171/JNCG/648}{doi:10.4171/JNCG/648}.
Inspected source: \href{https://arxiv.org/abs/2504.12730v2}{arXiv:2504.12730v2},
Remark~3.2, p.~11; Theorem~3.7, pp.~15--16.

\bibitem{CG} T.~Crisp and D.~Gow, Contractible subgraphs and Morita
equivalence of graph $C^*$-algebras, \emph{Proc. Amer. Math. Soc.}
\textbf{134} (2006), 2003--2013.
\href{https://doi.org/10.1090/S0002-9939-06-08216-5}{doi:10.1090/S0002-9939-06-08216-5}.
Inspected source: \href{https://arxiv.org/abs/math/0404542v1}{arXiv:math/0404542v1},
29 April 2004, Theorem~3.1.

\bibitem{GS} J.~Gabe and G.~Szab\'o, The dynamical Kirchberg--Phillips
theorem, \emph{Acta Math.} \textbf{232} (2024), 1--77.
\href{https://doi.org/10.4310/ACTA.2024.v232.n1.a1}{doi:10.4310/ACTA.2024.v232.n1.a1}.
Inspected sources: \href{https://arxiv.org/abs/2205.04933v1}{arXiv:2205.04933v1},
10 May 2022, Definition~1.15, Proposition~3.14, Theorem~5.8;
\href{https://arxiv.org/abs/2205.04933v3}{v3}, Proposition~3.15, Theorem~5.8.

\bibitem{HPP} A.~Hopenwasser, J.~R.~Peters and S.~C.~Power,
Subalgebras of graph $C^*$-algebras, \emph{New York J. Math.}
\textbf{11} (2005), 351--386, Theorem~5.2, p.~360.
\href{https://nyjm.albany.edu/j/2005/11-18.html}{Publisher article}.
The paper's source/range convention is reversed from ours.

\bibitem{Izumi} M.~Izumi, Finite group actions on $C^*$-algebras with
the Rohlin property, I, \emph{Duke Math. J.} \textbf{122} (2004), 233--280.
\href{https://doi.org/10.1215/S0012-7094-04-12221-3}{doi:10.1215/S0012-7094-04-12221-3}.
Inspected source: \href{https://www.math.kyoto-u.ac.jp/preprint/2002/4.ps}{Kyoto 2002 preprint no.~4},
Definition~3.6 and Remark~3.7, p.~13; Proposition~5.6(2), pp.~28--29.

\bibitem{KPR} A.~Kumjian, D.~Pask and I.~Raeburn,
Cuntz--Krieger algebras of directed graphs, \emph{Pacific J. Math.}
\textbf{184} (1998), 161--174.
\href{https://doi.org/10.2140/pjm.1998.184.161}{doi:10.2140/pjm.1998.184.161}.

\bibitem{Matsumoto} K.~Matsumoto, On automorphisms of $C^*$-algebras
associated with subshifts, \emph{J. Operator Theory} \textbf{44} (2000),
91--112, Lemma~5.14, p.~109, and Proposition~6.4, p.~111.
\href{https://jot.theta.ro/jot/archive/2000-044-001/2000-044-001-005.pdf}{Publisher PDF}.
\end{thebibliography}
\end{document}
